\documentclass[11pt, a4paper]{report}

\usepackage[utf8]{inputenc}
\usepackage[T1]{fontenc}
\usepackage[french]{babel}

% Appel des packages dans le dossier preambule
\def\packagepath{../../../../preambule}   % path du package princial

\usepackage{\packagepath/page_layout} % <--- Nouveau
\usepackage{\packagepath/config_info}
\usepackage{\packagepath/cadres_cours}
\usepackage{\packagepath/zones_reponse}
\usepackage{\packagepath/config_ds}
\usepackage{\packagepath/config_maths}

% --- PILOTAGE ---
\def\visibleornot{visible}    % visible or invisible


\def\documentpath{./Sources_Latex}
\graphicspath{{./Images/}}


\def\ptsexoA{0}

\begin{document}

%%%%%%%%%%%%%%%%%%%%%%%%%%%
\def\level{BTS}  % Classe à droite
\def\course{CIEL 1}
\def\chaptername{Intégration} % Titre au centre
\def\docname{C107 -- TP3}        % Identifiant à gauche

\pagestyle{DS_FP}
%%%%%%%%%%%%%%%%%%%%%%%%%%%%%%%%%

\NomPrenomNote{}

%%=============================================================
\section*{PARTIE A: Détermination expérimentale du paramètre $c$}

On cherche à modéliser un signal par la fonction suivante :  

\[f(x) =  x\cdot\e^{-c\cdot x} \qquad \text{ (où $c>0$ est le paramètre à déterminer).}\]

Le signal doit respecter la contrainte physique suivante:

\begin{itemize}
    \item le signal doit présenter un maximum à l'instant $x=0,5$ s.
\end{itemize}

\begin{enumerate}
    \item Tracer la courbe de $f$ pour différentes valeurs de $c$. 
    \item Identifier la valeur de $c$ qui satisfait les contraintes données. 
     \begin{blocvisible}[height=1]
        $c = 2$ satisfait les contraintes, car $f'(0,5) = 0$.
    \end{blocvisible}
\end{enumerate}


\section*{PARTIE B: Étude du signal identifié}

Utilisez maintenant la valeur trouvée ($c = 2$) pour analyser le comportement du signal $f(x) = x \cdot\e^{-2x}$.

\begin{enumerate}[resume]
    \item Tracer la courbe représentative de $f$ et observer son comportement. 
    \item Affichez la fonction dérivée $f'$  
    \item Déterminer le signe de $f'$ et en déduire les variations de $f$. 
     \begin{blocvisible}[height=4]
        \begin{minipage}{0.58\linewidth}
            $f'(x) = \e^{-2x} - 2x \e^{-2x} = \e^{-2x}(1 - 2x)$, qui est positif pour $x < 0,5$ et négatif pour $x > 0,5$.

        Par conséquent, $f$ est strictement croissante sur $\ff{0}{0,5}$ et strictement décroissante sur $\ff{0,5}{+\infty}$.
        \end{minipage}\hfill
        \begin{minipage}{0.4\linewidth}
            \begin{tikzpicture}
                \tkzTabInit[lgt=2,espcl=3.5]
                {$x$ / 1, $signe \, f'(x)$ / 1, $var \, f$ / 1}
                {$0$, $+\infty$}
                \tkzTabLine{}
                \tkzTabVar{}
            \end{tikzpicture}
        \end{minipage}
    \end{blocvisible}
    \item Placer un point $N\left(x_N; f(x_N)\right)$ sur la courbe et observer la valeur de $f(x_N)$ pour différentes positions de $N$.        
    \item Déplacez le point $N$ vers la droite (grandes valeurs de $x$) et observez la tendance de $f(N)$. 
    \item Conjecturez la limite de $f(x)$ lorsque $x$ tend vers $+\infty$.  
    
    \begin{blocvisible}[height=3]
        Lorsque $x$ tend vers $+\infty$, $f(x)$ tend vers 0, car la fonction exponentielle décroît rapidement.

        On peut également calculer la limite formellement :
        \[\lim_{x \to +\infty} x \e^{-2x} = \lim_{x \to +\infty} \frac{x}{\e^{2x}}=  0\]
    \end{blocvisible}
\end{enumerate}

\bigskip

\pagebreak
\thispagestyle{DS_LP}

\section*{PARTIE C: Intégration et Valeur Moyenne}

On souhaite évaluer la tension moyenne délivrée par le composant sur un intervalle $[a ; b]$.

\begin{enumerate}[resume]
    \item Créer les curseurs $a$ et $b$ avec les valeurs initiales $a=0$ et $b=2$.  
    \item Déterminer graphiquement l'aire sous la courbe de $f$ entre $a$ et $b$ avec la commande \verb|Aire = Intégrale(f, a, b)|.  
    \begin{blocvisible}[height=2]
        D'après Géogébra, l'aire sous la courbe de $f$ entre $a=0$ et $b=2$ est d'environ $0,38$ unités d'aires.
    \end{blocvisible}
        
    \item Observer comment l'aire change lorsque vous modifiez les valeurs de $a$ et $b$. 
    \begin{blocvisible}[height=3]
        Lorsque vous augmentez $b$, l'aire sous la courbe de $f$ entre $a$ et $b$ augmente, car vous incluez plus de la courbe. Inversement, lorsque vous augmentez $a$, l'aire diminue, car vous excluez une partie de la courbe.
    \end{blocvisible}

    \item A l'aide de Géogébra, déterminer l'expression d'une primitives $F$ de $f$. 
    \begin{blocvisible}[height=2]
        Une primitive de $f(x) = x \cdot\e^{-2x}$ est donnée par :
        \[F(x) = -\frac{1}{2}\cdot x \cdot\e^{-2x} - \frac{1}{4}\cdot \e^{-2x} + C= -\frac{(2x + 1)}{4}\e^{-2x} + C\]
    \end{blocvisible}

    \item Calculer l'aire sous la courbe de $f$ entre $a=0$ et $b=2$ à l'aide de l'intégrale. 
    \begin{blocvisible}[height=3]
        \[\int_0^2 f(x) \dx = F(2) - F(0) = \left(-\frac{1}{2} \cdot 2 \cdot \e^{-4} - \frac{1}{4} \e^{-4}\right) - \left(-\frac{1}{4}\right) = 0,38\]

        L'aire calculée à l'aide de l'intégrale est d'environ $0,38$ unités d'aires, ce qui correspond à l'aire trouvée graphiquement.
    \end{blocvisible}
    
    \item Déterminer la valeur moyenne $\mu$ de la fonction $f$ sur l'intervalle $[a ; b]$. 
    \begin{blocvisible}[height=3]
        \[\mu = \frac{1}{b-a} \cdot\int_a^b f(x) \, dx = \frac{1}{2-0} \cdot 0,38 = 0,19\]

        La valeur moyenne de la fonction $f$ sur l'intervalle $\ff{0}{2}$ est d'environ $0,19$ volts.
    \end{blocvisible}
    
    \item Créer un rectangle d'équivalence avec les points \verb|(a,0)|, \verb|(b,0)|,\verb|(b, moy)| et \verb|(a, moy)|. 
    
    Utilisez l'outil \verb|Polygone| pour le colorer. 

    \item En conservant $a=0$, déterminer la valeur de $b$ pour laquelle la valeur moyenne $\mu$ de $f$ sur l'intervalle $\ff{a}{b}$ est égale à $0,1$ volts (à $10^{-2}$ près). 
    \begin{blocvisible}[height=3]
        A l'aide de Géogébra, on fixera $a=0$ et on fera varier $b$ jusqu'à ce que la valeur moyenne $\mu$ soit égale à $0,1$ volts.

        On obtient une valeur de $b$ d'environ $1,25$.
    \end{blocvisible}
\end{enumerate}

\end{document}

